\documentclass[10pt]{article}
\usepackage[margin=20mm]{geometry}
\usepackage{amsmath,amssymb,mathtools,fontspec,unicode-math,xurl,hyperref}
\IfFontExistsTF{DejaVu Serif}{\setmainfont{DejaVu Serif}}{\setmainfont{Latin Modern Roman}}
\setmathfont{latinmodern-math.otf}
\setlength{\parindent}{0pt}\setlength{\parskip}{6pt}\sloppy
\title{Scientific Review of
Circle Bundles with PSC
over Large Manifolds}
\author{AI-assisted technical review}
\date{}
\begin{document}
\maketitle
Abstract: The Donaldson-divisor construction gives a compelling mechanism for circle bundles with positive scalar curvature over large bases. Theorem A survives the local corrections established here. One Major problem concerns the four-dimensional skeletal-compression dependency used to finish Theorem B: a decisive homotopy-group identification in the cited proof is false in the required dimension. This is a confirmed gap in the cited argument, not a counterexample to Theorem B. Three Minor findings concern the boundary-to-spine proof, an impossible alternative in Example 10, and a width constant. Explicit witnesses, repairs where available, and the limits of each conclusion are supplied below.\par
Identity and reviewed version. M25-0394: Aditya Kumar and Balarka Sen, Circle Bundles with PSC over Large Manifolds, Geometric and Functional Analysis 35 (2025), 1400–1423, DOI `10.1007/s00039-025-00723-z`. The packet’s exact scientific source is arXiv:2412.01047v5, 29 October 2025, 24 pages. Unless otherwise stated, page numbers are physical pages of that PDF, agreeing with its printed pagination. A journal PDF retrieved from the German National Library was separately retained and its identity and selected corresponding passages checked. This limited comparison is described below; full equivalence of the two versions is not assumed.\par
\section*{Scope, contribution, and overall assessment}
All 24 pages of the selected v5 were read, including definitions, proof details, examples, and references. The audit traced the symplectic divisor, its normal circle bundle, the complement’s handle indices, the fundamental-group identifications, rational essentiality, spin covers, the low-dimensional obstruction, and the subsequent examples. No substantive previous report or calculation for this paper was recoverable; the preserved sources and intake records were reused.\par
For Theorem A, a handle of index \(j\le n+1\) in a \((2n+2)\)-dimensional complement induces a surgery on its \((2n+1)\)-dimensional boundary of codimension \(2n+2-j\ge n+1\ge3\). This meets the PSC surgery range, including \(n=2\). The boundary group agrees with that of the complement by the relative handle range. Its lift admits a Lipschitz uniformly cobounded map to the lifted spine; Problem Problem 2 (\#p:spine) gives a complete replacement for the deficient diameter argument.\par
The symplectically aspherical and amenable case uses a rationally essential divisor and the exact amenability hypothesis of Dranishnikov’s Theorem 1.6 [Dra11]. For the product case, the projection has positive degree, and the even value of \(k\) makes \(c_1(\widehat Z)=\widehat j^{\,*}c_1(\widehat M)-k\widehat j^{\,*}[\widehat\omega]\) even on the selected spin cover. Thus the stated finite-cover definition of enlargeability is respected. The Euler characteristic is a polynomial in \(k\) with nonzero leading term, so it supplies infinitely many distinct homeomorphism types along large even \(k\). The parity adjustment in Example 11 agrees with the twistor Chern-class formula and successive adjunction in Fine–Panov [FP]. The positive-degree map in Example 16 has a proper lift, since it induces a fundamental-group isomorphism; the nonzero locally finite fundamental class fits [Retained latex: \textbackslash{}cite[Theorem 5.4]\{Dra13\}]. The full-dimension product claim on p. 5 fits [Retained latex: \textbackslash{}cite[Theorem 5.7]\{Dra13\}]; no general additivity assertion for arbitrary metric spaces is needed.\par
The first-review assessment is Major: one core dependency needs a valid four-dimensional replacement. Three local findings are Minor, with explicit corrections that preserve the surrounding conclusions. No independent acceptance decision is made.\par
\section*{The four-dimensional obstruction}
Problem 1 (Major: the final dependency in Theorem B).  The last paragraph on p. 18 invokes [Retained latex: \textbackslash{}cite[Main Theorem and Remark 1.7]\{Bol09\}] to deform the classifying map of the spin four-manifold \(E\) into \(B\pi^{(2)}\). The proof of that compression in the cited paper uses a false relative-to-quotient homotopy identification when \(n=4\). The same invocation appears on journal p. 1417 (physical p. 18 of the separately acquired journal PDF).\par
\subsection*{What is needed and what was checked}
After finite covers, the manuscript has an oriented base with torsion-free three-manifold group and an oriented circle bundle. The nontrivial-fiber case is covered by the hypotheses of He’s Theorem 1.9 [He]: there is an aspherical prime summand and the total dimension is four. In the remaining case the fiber is null in \(\pi_1(E)\), so \(\pi_1(E)\cong\pi_1(M)\) and \(\operatorname{cd}\pi\le3\). The splitting \(TE\cong p^*TM\oplus\mathbb R\) does make \(E\) spin. The Baum–Connes result for three-manifold groups is the relevant assembly input [MOP]. Its use here requires the integral real assembly conclusion, not merely rational complex injectivity; the full Baum–Connes isomorphism, with the standard real/complex comparison, is the appropriate route.\par
Lemma 3.2 prevents a classifying map from landing in the two-skeleton. It does not, by itself, turn a nonequivariant uniformly cobounded map \(\widetilde E\to K^2\) into such a classifying map. Thus a bound on \(\dim_{\mathrm{mc}}\widetilde E\) alone cannot replace the compression assertion.\par
\subsection*{A decisive calculation inside the stated group range}
Take \(n=4\), \(\pi=\mathbb Z^3\), and the standard cubical model \(B\pi=T^3\). This is a torsion-free group of cohomological dimension \(3=n-1\); it satisfies the assembly hypothesis. Let \[X=\widetilde{B\pi}=\mathbb R^3,\qquad A=\widetilde{B\pi^{(2)}}.\] Here \(A\) is the two-skeleton of the unit cubulation. Choose a closed unit cube \(Q\) and its center \(c\). Every point of \(A\) lies outside the interior of \(Q\). Radial projection from \(c\), using the cubical norm, defines \[r:A\longrightarrow\partial Q\cong S^2,\qquad
 r|_{\partial Q}=\operatorname{id}.\] For example, after centering \(Q\) at the origin with side length one, \(r(x)=x/(2\|x\|_\infty)\). Consequently the inclusion \(i:S^2\hookrightarrow A\) has a retraction and \(i_*:\pi_3(S^2)\to\pi_3(A)\) is injective. The Hopf class \(\eta\in\pi_3(S^2)\cong\mathbb Z\) therefore gives an infinite-order element \(i_*\eta\).\par
As \(X\) is contractible, the exact sequence of the pair yields \[\pi_4(X,A)\cong\pi_3(A).\] On the other hand, collapsing the two-skeleton identifies \[X/A\cong\bigvee_{\gamma\in\mathbb Z^3}S^3,\qquad
 \pi_4(X/A)\cong\bigoplus_{\gamma\in\mathbb Z^3}\mathbb Z/2.\] There are no Whitehead cross-terms in degree four for a wedge of three-spheres. The relative class represented inside \((Q,\partial Q)\) by \(\eta\) maps to the suspended Hopf class in its \(S^3\) summand. In particular, twice that relative class is nonzero, but its image in \(\pi_4(X/A)\) is zero.\par
This directly disproves the identification \[\pi_n(\widetilde{B\pi},\widetilde{B\pi}^{(n-2)})
 \ \cong\
 \pi_n(\widetilde{B\pi}/\widetilde{B\pi}^{(n-2)})
 \tag{*}\] used on printed p. 24 of [Bol09], together with the assertion there that the left side is a free \(\mathbb Z_2[\pi]\)-module. All hypotheses of that purely topological identification are met. No PSC assumption enters the asserted identification. This calculation is not presented as a PSC four-manifold violating Theorem B.\par
The issue is exactly an unstable-dimensional one: the map \(\pi_3(S^2)\to\pi_4(S^3)\) loses twice the Hopf class. Detecting the suspended class in \(KO\) cannot recover this lost information without an additional argument. The subsequent bordism calculation on pp. 25–26 of [Bol09] does not establish that missing injectivity.\par
\subsection*{Counterreadings, impact, and repair limits}
The main theorem of [Bol09] is stated in a range including \(n=4\), and its proof explicitly concludes skeletal compression. Thus the manuscript has not simply misquoted a dimension restriction. The problem is the invalid step in the proof supporting its decisive citation. The calculation above establishes that gap; it does not prove the stated compression theorem false.\par
Replacing the citation by Bolotov–Dranishnikov [Retained latex: \textbackslash{}cite[Lemma 4.1]\{BD\}] is not an established repair. Dranishnikov explicitly records a gap in that lemma in [Retained latex: \textbackslash{}cite[Remark 1.2]\{Dra21\}]; his treatment excludes almost-spin four-manifolds. The later argument in [DH] also gives its relevant strong-inessentiality theorem for \(n>4\). These observations corroborate the need to check the dimension, but the independent calculation above, rather than a general literature inference, is the decisive evidence.\par
The gap affects the null-fiber branch of Theorem B and therefore the claimed answer in all dimensions. It does not affect the high-dimensional existence construction in Theorem A or He’s nontrivial-fiber branch. A valid proof specialized to these circle bundles, a direct index obstruction, or another verified four-dimensional result could repair the argument. Such a replacement has not been completed here. Increasing a dimension bound to five would exclude the case that is needed. The current finding is a Major proof-dependency gap, with the truth of Theorem B left open by this review.\par
\section*{The boundary-to-spine argument}
Problem 2 (Minor: retraction and intrinsic diameter).  The introductory deformation-retraction claim on p. 5 is too strong. In Lemma 3.1, pp. 12–13, the assertion that a finite family of inverse images of open stars has a finite maximum intrinsic diameter is unjustified and can fail. The lemma’s conclusion has the complete same-hypothesis repair below.\par
\subsection*{The boundary need not retract to the spine}
Let \(M=\mathbb{CP}^3\) with its integral hyperplane symplectic class and let \(Z\subset\mathbb{CP}^3\) be a smooth complex hypersurface of sufficiently large degree \(k\). This is an allowed Donaldson divisor with \(n=2\). Its canonical circle bundle has Euler class \(kH|_Z\). Lefschetz gives \(\pi_1(Z)=0\) and a surjection \(\pi_2(Z)\to\pi_2(\mathbb{CP}^3)\cong\mathbb Z\), so the bundle sequence gives \[\pi_1(E)\cong\mathbb Z/k.\] Thus \(\widetilde E\) is a closed oriented five-manifold and \(H_5(\widetilde E;\mathbb Z)\cong\mathbb Z\). A three-dimensional simplicial complex has zero fifth homology. There can be no deformation retraction of \(\widetilde E\) onto the proposed three-dimensional spine. The complement, rather than its boundary, retracts; composing boundary inclusion with that retraction is sufficient. Indeed \(\widetilde E\) is compact in this example, so its macroscopic dimension is zero. The example refutes only the retraction assertion.\par
\subsection*{A counterexample to the stated intrinsic-diameter step}
Use \[W=S^1\times D^3,\qquad E=S^1\times S^2,\qquad K=S^1\times\{0\}.\] Both hypotheses of Lemma 3.1 hold and \(E\to W\) induces an isomorphism on \(\pi_1\). For \(x\in[0,1]\), let a degree-one lift of a circle map be \[H(x)=
 \begin{cases}
  3x,&0\le x\le1/3,\\
  2-3x,&1/3\le x\le2/3,\\
  3x-2,&2/3\le x\le1,
 \end{cases}
 \qquad H(x+1)=H(x)+1.\] Write \(h_t([x])=[(1-t)x+tH(x)]\). Then \[R_t(s,z)=\bigl(h_{t|z|}(s),(1-t)z\bigr)\] is a strong deformation retraction from \(W\) onto \(K\). Its terminal map, restricted to \(E\), is \(f(s,z)=h_1(s)\). The boundary map is simplicial after finite triangulations.\par
Choose an open star corresponding to the arc \(B=(1/4,3/8)\) of the circle. Its inverse image is the product of \(S^2\) with three disjoint intervals: \[(1/12,1/8),\qquad(13/24,7/12),\qquad(3/4,19/24).\] There is no path inside \(U=f^{-1}(B)\) joining different components. Its intrinsic path distance between those components is infinite. Finiteness of the index set of stars therefore does not supply the finite constant \(C\) asserted on p. 13. The exact inverse-interval arithmetic was independently executed in the retained reviewer script.\par
\subsection*{Complete repair with unchanged hypotheses}
Replace \(f\) by its finite simplicial approximation, which is Lipschitz. Because \(f_*\) is an isomorphism, its lift \(\widetilde f:\widetilde E\to\widetilde K\) is equivariant for the same deck group \(\Gamma\). Choose compact sets \(D_E,D_K\) whose translates cover the two universal covers. Both \(\widetilde f(D_E)\) and \(D_K\) are compact. Properness of the deck action implies that \[F=\{\gamma\in\Gamma:
       \gamma\widetilde f(D_E)\cap D_K\ne\varnothing\}\] is finite. It follows that \[\widetilde f^{-1}(D_K)\subseteq\bigcup_{\gamma\in F}\gamma D_E.\] The right side has finite diameter in the lifted Riemannian metric. For every target point, a deck translate moves it into \(D_K\); equivariance and the isometric deck action give the same diameter bound for every fiber. The lift is Lipschitz, so \(\dim_{\mathrm{MC}}\widetilde E\le k\) follows.\par
This avoids any connectedness or intrinsic-diameter assumption on inverse images of stars. The introductory sentence should describe this map rather than a retraction of the boundary. The correction is Minor because a full proof of the same lemma and the same Theorem A bound is supplied, without new hypotheses or a restricted family of examples. The analogous intrinsic-diameter passage is also present on journal p. 1411.\par
\section*{Two local quantitative corrections}
Problem 3 (Minor: an impossible branch in Example 10).  The final alternative in Example 10, p. 19, proposes that when \(m=2\) the fiber homomorphism can have nonzero kernel, producing a finite cyclic extension and dimensions \(2n\). For the canonical bundle actually constructed here, that alternative cannot occur. It also appears on journal p. 1417.\par
Let \(j:Z^{2m}\hookrightarrow T^{2n}\) be any of the stated iterated divisors. The restricted symplectic class is \(j^*[\omega]\). For every \(a:S^2\to Z\), \[\langle j^*[\omega],a_*[S^2]\rangle
 =\langle[\omega],(j\circ a)_*[S^2]\rangle=0,\] because \(\pi_2(T^{2n})=0\). The canonical bundle in the final divisor step has Euler class \(e=k\,j^*[\omega]\) (with the relevant final positive integer \(k\)). The connecting homomorphism in its homotopy sequence is Euler evaluation: \[\partial:\pi_2(Z)\longrightarrow\pi_1(S^1)\cong\mathbb Z,\qquad
 \partial[a]=\langle e,a_*[S^2]\rangle=0.\] This is valid also for \(m=2\), even though \(\pi_2(Z)\) may then be nonzero. Hence the actual group is always a central extension \[1\longrightarrow\mathbb Z\longrightarrow\pi_1(E)
   \longrightarrow\mathbb Z^{2n}\longrightarrow1.\] It is torsion-free and polycyclic of Hirsch length \(2n+1\), so the same polycyclic-group formulas invoked by the paper give \(\mathop{\mathrm{asdim}}\pi_1(E)=\mathop{\mathrm{vcd}}\pi_1(E)=2n+1\). The manuscript’s calculation in Example 15, p. 22, supplies a useful internal cross-check: it also makes the pulled-back Euler class vanish.\par
This finding does not declare a conditional statement with an impossible antecedent logically false. It corrects the claim that this alternative is available in the described family. Remove that alternative and use the infinite-cyclic extension uniformly. The intended contrast between large group dimension and small total-space macroscopic dimension remains, and is stronger in the \(m=2\) case.\par
Problem 4 (Minor: the round-sphere width constant).  On p. 3 the projection of \(\widetilde M\times S_r^2\) is used to claim \(\mathrm{UW}_{n-2}\le2r\) for the product Riemannian metric. The fiber’s intrinsic diameter is \(\pi r\), not \(2r\).\par
The unqualified statement on pp. 2–3 includes \(n=2\) and \(M\) a point, a closed zero-dimensional manifold. Every continuous map from the connected sphere to a zero-dimensional simplicial complex is constant. Consequently \[\mathrm{UW}_0(S_r^2)=\mathop{\mathrm{diam}}(S_r^2)=\pi r>2r.\] This gives an exact counterexample to the literal displayed bound. In every dimension, the projection proves the corrected upper bound \(\mathrm{UW}_{n-2}\le\pi r\). The scalar-curvature scaling argument and the macroscopic-dimension conclusion are unchanged. Even if an implicit \(n>2\) restriction were intended, the stated projection argument still supplies \(\pi r\).\par
\section*{Counterevidence, dependencies, and completion record}
The Steenrod-square line preceding (3.4) on p. 17 deserves a clearer justification, but it does not yield an additional theorem-level defect. Degree one by itself does not imply \(\mathrm{Sq}^1z=0\). Here the oriented rank-two bundle has \(\mathrm{Sq}^1u=w_1u=0\), while \(\mathrm{Sq}^2z=0\) by instability, so the unwanted Cartan terms vanish. Moreover the actual \(z\) obtained by Thom inversion of \(F^*\eta\) lifts integrally, because \(\eta\) and the oriented Thom class lift integrally. This also makes \(\mathrm{Sq}^1z=0\) for the specific class. An arbitrary degree-one class with nonzero square would not satisfy these provenance conditions and is not used as a counterexample.\par
The finite-cover reduction in the proof of Theorem B is consistent with Kurosh’s subgroup theorem: the kernel of the map to the product of the finite prime factors misses their nontrivial torsion. The term abelianization used earlier for the free-product-to-product map should mean the canonical product homomorphism when a factor is nonabelian. This is a terminology correction, not a separate obstruction. Similarly, the dimension label for the punctured complement in the proof of Theorem A should be \(2n+2\).\par
The primary-reference check was recorded by actual scope: Bolotov [Bol09], all seven pages, with the central four proof pages visually inspected; [BD], pp. 3 and 7–8; [Dra21], pp. 3 and 12–13; [DH], pp. 3, 9 and 12; [He], pp. 4–5 and 18–21; [MOP], physical pp. 5–7; [Dra11], pp. 3 and 16–17 of the retained preprint; [Dra13], pp. 18–19; and [FP], all four pages. The earlier Thom calculation in [Bol06], pp. 4–5, was also checked. These are bounded dependency checks; downloading or parsing the remaining pages is not counted as scientific reading. The attempted publisher access to Donaldson’s original article encountered a security check; no bypass was attempted. Its standard existence and Morse results are used as stated by the manuscript, with their application and the manuscript’s deductions checked, rather than claiming a new full audit of that external proof.\par
The separate journal PDF was checked at physical pp. 1, 12–13 and 18–19. It confirms the identity and the corresponding passages identified above; it is not labeled a full journal-version review. The frozen source manifest and selected v5 remain the principal input. No missing code is treated as a mathematical defect; no author code was executed, and no unseen supplement is credited.\par
Issue | Severity | Actual impact\par
P1 | Major | Theorem B’s null-fiber branch needs a valid four-dimensional proof dependency; no theorem counterexample claimed.\par
P2 | Minor | Boundary retraction and star-diameter argument corrected; the same lemma and Theorem A bound retained.\par
P3 | Minor | Example 10 has only the infinite-cyclic extension branch.\par
P4 | Minor | Replace \(2r\) by \(\pi r\); macroscopic conclusion retained.\par
The proofs and exact arithmetic are the work of the current scientific conversation’s root assistant. The controller’s supplied observation of the official 6Pro product route is retained separately; backend equivalence is not independently asserted.\par
This English report was actually compiled with the authorized Tectonic 0.17.0 runtime. Per-paper cache, payloads, commands, exit codes, hashes, logs, PDF parsing and each page’s native visual-inspection evidence are retained with the delivery. This is first-review recovery only: `scientific\_decision=pending`, empty human reviews, and independent scientific acceptance deferred to a separate later review.\par
99 D. Bolotov, About the macroscopic dimension of certain PSC-manifolds, Algebr. Geom. Topol. 9 (2009), 21–27. https://msp.org/agt/2009/9-1/agt-v9-n1-p02-p.pdf. D. Bolotov and A. Dranishnikov, On Gromov’s scalar curvature conjecture, arXiv:0901.4503v1 (2009). A. Dranishnikov, Positive scalar curvature and strongly inessential manifolds, arXiv:2105.07528v1 (2021). A. Dranishnikov and S. Howladar, On Gromov’s conjecture for right-angled Artin groups, arXiv:2311.00573v1 (2023). S. He, Twisted \(S^1\) stability and PSC obstruction on fiber bundles, arXiv:2303.12614v4 (2025). M. Matthey, H. Oyono-Oyono and W. Pitsch, Homotopy invariance of higher signatures and 3-manifold groups, Bull. Soc. Math. France 136 (2008), 1–25. A. Dranishnikov, On macroscopic dimension of rationally essential manifolds, retained source arXiv:1005.0424v1 (2010). A. Dranishnikov, On macroscopic dimension of universal coverings of closed manifolds, arXiv:1307.1029v1 (2013). J. Fine and D. Panov, Symplectic domination, arXiv:1905.05671v1 (2019). D. V. Bolotov, Gromov’s macroscopic dimension conjecture, Algebr. Geom. Topol. 6 (2006), 1669–1676; retained journal-text copy arXiv:0904.4886v1.\par
\end{document}
